\documentclass[11pt]{article}
\usepackage[margin=1in]{geometry}
\usepackage{amsmath,amssymb,amsthm,mathtools}
\usepackage{enumitem}
\usepackage{booktabs}
\usepackage{hyperref}
\hypersetup{colorlinks=true,linkcolor=blue,citecolor=blue,urlcolor=blue}

\newtheorem{theorem}{Theorem}[section]
\newtheorem{lemma}[theorem]{Lemma}
\newtheorem{proposition}[theorem]{Proposition}
\newtheorem{corollary}[theorem]{Corollary}
\newtheorem{definition}[theorem]{Definition}
\newtheorem{remark}[theorem]{Remark}
\newtheorem{example}[theorem]{Example}

\newcommand{\C}{\mathbb C}
\newcommand{\R}{\mathbb R}
\newcommand{\Ran}{\operatorname{Ran}}
\newcommand{\Ker}{\operatorname{Ker}}
\newcommand{\tr}{\operatorname{tr}}
\newcommand{\Capacity}{\operatorname{Cap}}
\providecommand{\Cost}{\operatorname{Cost}}
\newcommand{\XiC}{\Xi_C(D\mid L)}
\newcommand{\preceqq}{\preceq}
\newcommand{\eps}{\varepsilon}

\title{Adequacy Residuals and Blind-Spot Currency\\
\large A Companion Paper to Formed-Layer Membrane Theory}
\author{Working draft v0.1}
\date{May 2026}

\begin{document}
\maketitle

\begin{abstract}
Formed-layer membrane theory distinguishes a layer's native probe family from a larger family of probes that would dissolve the layer if they were cheaply addressable.  This note isolates the object that measures the difference.  Given an audit energy $C\succeq 0$, a native probe family $L$, and a layer-dissolving probe family $D$, the \emph{adequacy residual}
\[
\Xi_C(D\mid L)=K_{DD}-K_{DL}K_{LL}^{\dagger}K_{LD}
\]
is the Schur complement of dissolving currency after conditioning on native currency.  Equivalently,
\[
\Xi_C(D\mid L)=T_D(I-P_L)T_D^*,
\]
so it is the component of the dissolving readout invisible to native probes in the energy-dual geometry.  The residual is the unique optimal explanation defect for factorizations $D\approx AL$, vanishes exactly when $\ker L_0\subseteq \ker D_0$ on the legal energy quotient, decreases under lawful strict native-family extension, and obeys data-processing only through audited bridges.  Thus $\Xi_C(D\mid L)$ is the blind-spot currency of a formed layer: what the layer-breaking probes can see that the layer's official native probes cannot.
\end{abstract}

\section{Purpose and scope}

This companion note extracts the adequacy residual from the larger formed-layer membrane theory.  The main membrane paper studies predictive native probe-family currency matrices
\[
K_{C,\Gamma}(L)=L_\Gamma C_\Gamma^\dagger L_\Gamma^*,
\]
critical-pair completion, exact-six records, and root-composite obligations.  The present note focuses on one object:
\[
\Xi_C(D\mid L),
\]
where $L$ is the native probe family of a formed closure/layer and $D$ is a family of layer-dissolving probes.

The guiding interpretation is:
\[
\boxed{\Xi_C(D\mid L)=\text{what the dissolving probes can see that the native probes cannot.}}
\]

This note does not claim that raw substrates have no needles.  It applies only after the carrier, audit, null-mode quotient, native probe family, and dissolving probe family have been declared.  Non-closed artifacts and undeclared probes are outside scope.

\section{Standing hypotheses: the legal energy quotient}

Let $E$ be a finite-dimensional real or complex Hilbert space and let
\[
C=C^*\succeq 0
\]
be an audit energy.  The legal energy quotient is represented by
\[
E_C=(\ker C)^\perp,
\qquad
C_0=C|_{E_C}>0.
\]
All finite-currency formulas in this note are read on $E_C$.  A probe family $F:E\to Y$ is \emph{null-legal} if
\[
F|_{\ker C}=0.
\]
For null-legal $F$ we write $F_0=F|_{E_C}$.

\begin{remark}[Null-mode warning]
If a probe sees a legal zero-energy direction, the variational capacity is infinite even if a pseudoinverse expression appears finite.  Null-mode legality is therefore not a technicality; it is part of the adequacy record.
\end{remark}

\section{Currency and minimum spend}

Given a null-legal probe family $F:E\to Y$, define its currency matrix
\[
K_F=F_0C_0^{-1}F_0^*=FC^\dagger F^*.
\]
For a recombination $y\in Y$,
\[
\Capacity_C(y^*F)=y^*K_Fy.
\]
Thus $K_F$ is the compliance, or response-per-audit-dollar, matrix.

\begin{theorem}[Minimum-spend duality]
For $z\in Y$ define
\[
\Cost_C(z;F)=\inf\{\langle u,Cu\rangle:\ Fu=z\}.
\]
Then
\[
\Cost_C(z;F)=
\begin{cases}
 z^*K_F^\dagger z,& z\in \Ran F_0=\Ran K_F,\\
 +\infty,& z\notin\Ran F_0.
\end{cases}
\]
\end{theorem}

\begin{proof}
On $E_C$, set $v=C_0^{1/2}u$.  The constraint becomes $F_0C_0^{-1/2}v=z$.  The minimum norm solution has squared norm $z^*(TT^*)^\dagger z$ for $T=F_0C_0^{-1/2}$.  Since $TT^*=K_F$, the result follows.  If $z$ is outside the range there is no feasible $u$.
\end{proof}

\section{Definition of the adequacy residual}

Let
\[
L:E\to Y
\]
be the native probe family and let
\[
D:E\to Z
\]
be a layer-dissolving probe family.  Assume both are null-legal.  Define the block currency matrices
\[
K_{LL}=L_0C_0^{-1}L_0^*,
\qquad
K_{DL}=D_0C_0^{-1}L_0^*,
\]
\[
K_{LD}=L_0C_0^{-1}D_0^*,
\qquad
K_{DD}=D_0C_0^{-1}D_0^*.
\]

\begin{definition}[Adequacy residual]
The adequacy residual of $D$ relative to $L$ is
\[
\boxed{\Xi_C(D\mid L)=K_{DD}-K_{DL}K_{LL}^\dagger K_{LD}.}
\]
\end{definition}

It is the Schur complement of the native block inside the joint currency matrix
\[
\begin{pmatrix}
K_{LL}&K_{LD}\\
K_{DL}&K_{DD}
\end{pmatrix}.
\]

\section{Projection theorem}

Set
\[
T_L=L_0C_0^{-1/2},
\qquad
T_D=D_0C_0^{-1/2}.
\]
Let $P_L$ be the orthogonal projection onto $\Ran(T_L^*)\subseteq E_C$.

\begin{theorem}[Projection identity]
\[
\boxed{\Xi_C(D\mid L)=T_D(I-P_L)T_D^*\succeq 0.}
\]
\end{theorem}

\begin{proof}
Since $K_{LL}=T_LT_L^*$, the operator $T_L^*K_{LL}^\dagger T_L$ is the orthogonal projection onto $\Ran(T_L^*)$.  Therefore
\[
K_{DL}K_{LL}^\dagger K_{LD}=T_DT_L^*K_{LL}^\dagger T_LT_D^*=T_DP_LT_D^*.
\]
Subtracting from $K_{DD}=T_DT_D^*$ gives the identity.
\end{proof}

\begin{remark}
The theorem says that $\Xi_C(D\mid L)$ is the dissolving readout after removing the part explained by native probes in the energy-dual geometry.
\end{remark}

\section{Optimal residual and exact adequacy}

Define the optimal native explanation map
\[
A_*=K_{DL}K_{LL}^\dagger:Y\to Z.
\]

\begin{theorem}[Optimal residual identity]
For every linear map $A:Y\to Z$,
\[
\boxed{(D_0-AL_0)C_0^{-1}(D_0-AL_0)^*
=\Xi_C(D\mid L)+(A-A_*)K_{LL}(A-A_*)^*.}
\]
Consequently, $\Xi_C(D\mid L)$ is the Loewner-minimal residual currency among all attempts to factor $D$ through $L$.
\end{theorem}

\begin{proof}
Expand the left side:
\[
K_{DD}-AK_{LD}-K_{DL}A^*+AK_{LL}A^*.
\]
Insert $A=A_*+B$ and use $A_*K_{LL}=K_{DL}$ on the support of $K_{LL}$, which follows from the Moore--Penrose identities.  The cross terms cancel and the displayed identity remains.
\end{proof}

\begin{theorem}[Exact adequacy]
The following are equivalent:
\begin{enumerate}[label=(\roman*)]
\item $\Xi_C(D\mid L)=0$;
\item $D_0=A_*L_0$;
\item there exists $A$ such that $D_0=AL_0$;
\item $\ker L_0\subseteq \ker D_0$.
\end{enumerate}
\end{theorem}

\begin{proof}
By the projection theorem, $\Xi=0$ iff $T_D=T_DP_L$, i.e. $T_D^*$ lies in the range generated by $T_L^*$, equivalently $D_0$ factors through $L_0$.  The factorization criterion is precisely $\ker L_0\subseteq\ker D_0$.
\end{proof}

\section{Promotion from native membrane to dissolving membrane}

Suppose the native family has budget
\[
K_{LL}\preceq \Theta_Y
\]
and the adequacy residual has budget
\[
\Xi_C(D\mid L)\preceq \Xi_Z.
\]

\begin{theorem}[Adequacy promotion]
\[
\boxed{K_{DD}\preceq A_*\Theta_YA_*^*+\Xi_Z.}
\]
In particular, if $\Xi_C(D\mid L)=0$, native anti-localization transfers exactly to the dissolving family through $A_*$.
\end{theorem}

\begin{proof}
The Schur decomposition gives
\[
K_{DD}=A_*K_{LL}A_*^*+\Xi_C(D\mid L).
\]
Apply the two hypotheses.
\end{proof}

\section{Strict native extension}

Let $M:E\to W$ be an additional null-legal native probe family and set
\[
L^+=\begin{bmatrix}L\\M\end{bmatrix}.
\]
The conditional currency of $M$ given $L$ is
\[
K_{MM\mid L}=K_{MM}-K_{ML}K_{LL}^\dagger K_{LM},
\]
and similarly
\[
K_{DM\mid L}=K_{DM}-K_{DL}K_{LL}^\dagger K_{LM}.
\]

\begin{theorem}[Chain rule for adequacy residuals]
\[
\boxed{\Xi_C(D\mid L^+)=\Xi_C(D\mid L)-K_{DM\mid L}K_{MM\mid L}^\dagger K_{MD\mid L}.}
\]
In particular,
\[
\Xi_C(D\mid L^+)\preceq \Xi_C(D\mid L).
\]
\end{theorem}

\begin{proof}
This is the Schur complement chain rule applied to the joint currency block for $(L,M,D)$.  First condition on $L$, then take the Schur complement of the remaining $M$ block inside the residual joint currency of $(M,D)$.
\end{proof}

\begin{corollary}[Same-family saturation]
If $M=BL$ for some $B$, then
\[
\Xi_C(D\mid[L;M])=\Xi_C(D\mid L).
\]
\end{corollary}

\begin{proof}
If $M=BL$, then $K_{MM\mid L}=0$ and $K_{DM\mid L}=0$.
\end{proof}

\section{Data processing and transfer}

\begin{theorem}[Dissolving post-processing]
For a response map $F:Z\to Z'$,
\[
\boxed{\Xi_C(FD\mid L)=F\Xi_C(D\mid L)F^*.}
\]
\end{theorem}

\begin{proof}
Apply the definition to $FD$: all $D$-blocks are congruence-transformed by $F$.
\end{proof}

\begin{theorem}[Native coarsening]
For a native response map $A:Y\to Y'$,
\[
\boxed{\Xi_C(D\mid AL)\succeq \Xi_C(D\mid L).}
\]
If $A$ is a faithful reparameterization on $\Ran L_0$, equality holds.
\end{theorem}

\begin{proof}
The range of $(AL)_0^*$ is contained in the range of $L_0^*$.  Orthogonal projection onto a smaller explanation subspace leaves a larger residual.
\end{proof}

\begin{theorem}[Exact reducing bridge]
Suppose $U:E'_{C'}\to E_C$ is an isometric energy bridge with
\[
C'_0=U^*C_0U,
\quad
L'_0=AL_0U,
\quad
D'_0=BD_0U,
\]
where $A$ is faithful on the native range.  Then
\[
\boxed{\Xi_{C'}(D'\mid L')=B\Xi_C(D\mid L)B^*.}
\]
\end{theorem}

\begin{theorem}[Defective dissolving bridge]
If
\[
D'_0=BD_0U+R
\]
and the native bridge is exact, then for every $t>0$,
\[
\boxed{\Xi_{C'}(D'\mid L')
\preceq
(1+t)B\Xi_C(D\mid L)B^*+(1+t^{-1})\Xi_{C'}(R\mid L').}
\]
\end{theorem}

\begin{proof}
After projecting away the target native explanation space, use
\[
(X+Y)(X+Y)^*\preceq (1+t)XX^*+(1+t^{-1})YY^*.
\]
\end{proof}

\section{Obstruction and repair calculus}

For a budget $\Omega\succeq0$, define the adequacy defect
\[
\Delta_\Xi=\Xi_C(D\mid L)-\Omega.
\]

\begin{theorem}[Blind-spot witness]
\[
\Xi_C(D\mid L)\npreceq \Omega
\]
if and only if there exists $z\in Z$ such that
\[
\boxed{z^*(\Xi_C(D\mid L)-\Omega)z>0.}
\]
The maximum additive violation is $\lambda_{\max}(\Delta_\Xi)$.
\end{theorem}

\begin{proof}
This is the spectral characterization of the Loewner order.
\end{proof}

Lawful repair moves include:
\begin{enumerate}[label=(\alph*)]
\item strict native extension, which subtracts a positive Schur term;
\item bridge repair, which lowers the residual bridge currency;
\item null-mode repair, which quotients or annihilates illegal zero-energy readout;
\item budget repair, which weakens the claim but does not reduce the blind spot;
\item nonclaim/gating, which narrows scope rather than proving adequacy.
\end{enumerate}

\section{Predictive adequacy ladders}

At stage $j$, let
\[
\Xi_j=\Xi_{C_j}(D_j\mid L_j).
\]
If transported residuals satisfy
\[
\Xi_{j+1}\preceq (1+\eps_j)A_j\Xi_jA_j^*+E_j
\]
and the budgets obey the corresponding recurrence, adequacy propagates.  In a common response space, if
\[
\prod_j(1+\eps_j)<\infty,
\qquad
\sum_j E_j\preceq E_\infty,
\]
then predictive adequacy is maintained under a large enough budget.

A stronger collapse condition is
\[
\Xi_{j+1}\preceq (1-\alpha_j)\Xi_j+E_j,
\]
with
\[
\prod_j(1-\alpha_j)=0
\]
and controlled accumulated residuals.  This is the strict-extension form of blind-spot elimination.

\section{Examples and application templates}

\paragraph{Navier--Stokes.}
Let $L$ be energy/enstrophy/shell probes and let $D$ be pointwise concentration, vorticity, strain, or blow-up-rate probes.  Then $\Xi_C(D\mid L)$ measures the blind spot between classical energy accounting and blow-up control.  Energy/enstrophy alone do not adequately price pointwise blow-up probes in three dimensions; strict structure is required.

\paragraph{Root-composite/RH.}
Let $L$ be completed carrier-side anti-invariant readouts and let $D$ be off-fixed zero displacement readouts.  Then $\Xi_C(D\mid L)$ measures whether zero displacement is fully visible to the completed carrier currency.  The RH-facing obligation is not that $\Xi$ is automatically zero, but that explicit-formula domination and character-source coercivity earn the relevant adequacy and budget records.

\paragraph{Cryptography.}
Let $L$ be declared adversary/query families and $D$ unknown or extended attack vectors.  Then $\Xi$ is the security blind spot.

\paragraph{Immune systems.}
Let $L$ be known antigen recognition and $D$ novel pathogen readouts.  Then $\Xi$ is the immune blind spot.

\paragraph{Logic.}
Let $L$ be the declared proof/axiom readouts and $D$ broader formula-truth probes.  Then $\Xi$ measures the proof system's blind spot relative to the larger dissolving family.

\section{Nonclaims}

This paper does not claim that every raw substrate sees every dissolving probe.  It does not claim $\Xi=0$ generically.  It does not promote a public-shadow residual to an intrinsic residual without a bridge.  It does not say a finite-stage adequacy statement is predictive adequacy.  It states the exact residual, transfer, repair, and strict-extension laws that determine when such claims are valid.

\section{Conclusion}

The adequacy residual
\[
\Xi_C(D\mid L)
\]
is the blind-spot currency of a formed layer.  It is simultaneously a Schur complement, an optimal residual, a projection onto the part of dissolving readout invisible to native probes, and a monotone target of strict native-family extension.  It supplies the missing bridge between a native anti-localization membrane and a full layer-dissolving membrane.

\end{document}
